% !TeX root = ../../Book2.tex
% 纲要标题：### D.3 正性、GNS 与 Gelfand--Naimark 表示
% 纲要位置：工作记录/计划纲要.md，第二卷附录 D。

% 纲要要点（供目录核对）：
% * 正元素、正线性泛函和态（state）；归一化矩阵迹态的具体计算
% * 从态构造内积商空间并完备化，得到 Hilbert 空间上的循环 *-表示
% * 分别证明正锥、平方正性和次序性质，构造达到范数的态；证明 GNS、含幺和非含幺的 Gelfand--Naimark 表示定理及一般 \(*\)-同态的收缩性与单射等距性

\section{正性、GNS 与 Gelfand--Naimark 表示}
\label{b2:sec:D03}

线性泛函可以给代数元素赋予数值，却不必尊重乘法。要求它对平方取非负值，就能够把这些数值组织成内积，再从左乘构造表示。这是本节的主要构造。为让循环向量直接由单位元给出，本节先固定非零含幺 \(C^*\)-代数 \(A\)。

\subsection{正元素、正线性泛函与态}
\label{b2:subsec:D03:01}

\begin{definition}\label{b2:def:positive-element-state}
设 \(A\ne0\) 为含幺复 \(C^*\)-代数，\(a\in A\)。若 \(a=a^*\) 且其谱 \(\sigma_A(a)\subseteq[0,\infty)\)，则称 \(a\) 为\term[zheng-yuansu]{正元素}[positive element]，记 \(a\geq0\)，全体正元素组成 \(A_+\)。若复线性泛函 \(\varphi:A\to\mathbb C\) 对每个 \(a\in A_+\) 都满足 \(\varphi(a)\geq0\)，则称它为\term[zheng-xianxing-fanhan]{正线性泛函}[positive linear functional]。若正线性泛函 \(\varphi\) 还满足 \(\varphi(1_A)=1\)，则称它为 \(A\) 上的\term[tai]{态}[state]。
\symidx{\(A_+\)：\(C^*\)-代数的正元素锥}
\end{definition}

\begin{proposition}[正锥与连续函数演算]\label{b2:prop:cstar-positive-cone}
设 \(A\ne0\) 为含幺复 \(C^*\)-代数，\(A_+\) 为其正元素集合。则 \(A_+\) 是闭凸锥，且 \(A_+\cap(-A_+)=\{0\}\)。每个 \(p\in A_+\) 都有唯一正平方根 \(p^{1/2}\)。每个自伴元素 \(c\in A\) 都可写成
\[
c=c_+-c_-,\qquad c_\pm\in A_+,\quad c_+c_-=0,
\]
其中 \(c_\pm\) 由连续函数 \(t\mapsto\max(\pm t,0)\) 对 \(c\) 作函数演算得到。
\end{proposition}
\begin{proof}
对自伴元素 \(h\)，\cref{b2:thm:continuous-functional-calculus}把 \(C^*(1,h)\) 识别为实谱上的连续函数代数。因此 \(h\ge0\) 等价于谱上坐标函数非负，且
\[
h\ge0,\ \|h\|\le1
\quad\Longleftrightarrow\quad
h=h^*,\ \|h\|\le1,\ \|1-h\|\le1.
\]
右侧是三个闭凸集的交，故 \(A_+\) 与单位球的交闭且凸。将有界收敛序列或有限多个正元素同时缩放到单位球中，就得到正锥的闭性与凸性。非负实数倍仍为正，故它是闭凸锥。若 \(h\) 与 \(-h\) 都正，其谱只能是 \(\{0\}\)，自伴元素的范数等于谱半径，故 \(h=0\)。

对正元素的谱使用 \(t\mapsto\sqrt t\)，同一函数演算给出 \(p^{1/2}\)。对任意自伴元素 \(c\)，使用 \(t\mapsto\max(t,0)\) 和 \(t\mapsto\max(-t,0)\)，便得到所述分解及 \(c_+c_-=0\)。若正元素 \(q\) 满足 \(q^2=p\)，则 \(q\) 与 \(p\) 交换，也与 \(p^{1/2}\) 交换。在这两个元素生成的交换 \(C^*\)-代数中，非负函数的平方根逐点唯一，故 \(q=p^{1/2}\)。
\end{proof}

正平方根把谱上的非负性与代数中的平方联系起来。还须证明：即使 \(x\) 不是正规元素，\(x^*x\) 仍落在正锥中。

\begin{lemma}[平方元素的正性]\label{b2:lem:cstar-square-positive}
设 \(A\ne0\) 为含幺复 \(C^*\)-代数，则对每个 \(x\in A\)，都有 \(x^*x\in A_+\)。
\end{lemma}
\begin{proof}
对任意 \(u,v\in A\)，\(1-uv\) 可逆当且仅当 \(1-vu\) 可逆，因为一个方向的逆为
\[
(1-vu)^{-1}=1+v(1-uv)^{-1}u,
\]
另一方向交换 \(u,v\) 即得。对 \(\lambda\ne0\)，取 \(u=\lambda^{-1}y^*\)、\(v=y\)，便得
\[
\sigma_A(y^*y)\setminus\{0\}
=\sigma_A(yy^*)\setminus\{0\}.
\]
两者都是自伴元素，而可能多出的谱点 \(0\) 也属于 \([0,\infty)\)，所以其中一个为正当且仅当另一个为正。

写 \(x=a+ib\)，其中 \(a,b\) 自伴。由函数演算，\(a^2,b^2\ge0\)，再由正锥的加法封闭性得到 \(x^*x+xx^*=2a^2+2b^2\ge0\)。令 \(c=x^*x=c_+-c_-\)，取 \(y=xc_-\)。于是
\[
y^*y=c_-cc_-=-c_-^3,\qquad
yy^*=(y^*y+yy^*)+c_-^3\ge0.
\]
第二式对 \(y\) 使用同样的平方和等式，并使用 \(c_-^3\ge0\)。非零谱相同又使 \(y^*y\ge0\)；它的相反数也正，故 \(c_-^3=0\)。正元素 \(c_-\) 的谱因而只能是 \(\{0\}\)，于是 \(c_-=0\)。因此 \(x^*x\ge0\)。
\end{proof}

\begin{proposition}[正元素的次序性质]\label{b2:prop:cstar-positive-input}
设 \(A\ne0\) 为含幺复 \(C^*\)-代数。对自伴元素 \(s,t\in A\)，以 \(t-s\ge0\) 定义 \(s\le t\)。这给出自伴部分上的偏序；对任意 \(a,b\in A\)、正元素 \(p\in A_+\) 及自伴元素 \(h\in A\)，有
\[
0\le a^*a\le\|a\|^2 1,\qquad
b^*pb\ge0,\qquad
-\|h\|1\le h\le\|h\|1.
\]
\end{proposition}
\begin{proof}
正锥的加法封闭性和 \(A_+\cap(-A_+)=\{0\}\) 分别给出传递性和反对称性，\(0\in A_+\) 给出自反性。由上一个引理及 \(C^*\) 恒等式，\(a^*a\) 的谱包含于 \([0,\|a^*a\|]=[0,\|a\|^2]\)，故函数演算给出第一组界。又有
\[
b^*pb=(p^{1/2}b)^*(p^{1/2}b)\ge0.
\]
自伴元素 \(h\) 的谱位于 \([-\|h\|,\|h\|]\)，对 \(t\mapsto\|h\|\pm t\) 作函数演算便得到最后一组界。
\end{proof}

这些谱论与次序性质也见 Blackadar\cite[Propositions II.3.1.2--II.3.1.3; Corollary II.3.1.5; II.3.1.8]{Blackadar2017OperatorAlgebrasCorrected}。矩阵情形的正元素恰为半正定 Hermitian 矩阵，函数情形则是逐点非负的连续函数。

\begin{lemma}\label{b2:lem:positive-functional-cs}
设 \(A\ne0\) 为含幺复 \(C^*\)-代数，\(\varphi:A\rightarrow\mathbb C\) 为正线性泛函，则对任意 \(a,b\in A\)，有
\[
\varphi(a^*)=\overline{\varphi(a)},\qquad
\bigl|\varphi(a^*b)\bigr|^2\le\varphi(a^*a)\varphi(b^*b).
\]
它自动有界，且 \(\|\varphi\|=\varphi(1)\)。
\end{lemma}
\begin{proof}
对自伴 \(h\)，\cref{b2:prop:cstar-positive-input}使 \(\|h\|1\pm h\) 正，故 \(\varphi(h)\) 为实数。将任意 \(a\) 写成 \(h+ik\)、\(h,k\) 自伴，即得第一式。令 \([a,b]=\varphi(a^*b)\)，这是对第一变量共轭线性、第二变量线性的半正定 Hermitian 形式。

若 \([b,b]>0\)，代入 \(a-\lambda b\)、\(\lambda=[b,a]/[b,b]\)，由其平方值非负得到 \(\bigl|[a,b]\bigr|^2\le[a,a][b,b]\)。若 \([b,b]=0\)，非负式 \([a+zb,a+zb]\ge0\) 对任意复数 \(z\) 成立，其关于 \(z\) 的线性项只能为零，故 \([a,b]=0\)。这证明两种情形。

再把不等式用于 \(a,1\)，由 \(\varphi(a^*)=\overline{\varphi(a)}\) 及 \(0\le a^*a\le\|a\|^2 1\) 得
\[
\bigl|\varphi(a)\bigr|^2\le\varphi(a^*a)\varphi(1)
\le\|a\|^2\cdot\varphi(1)^2.
\]
所以 \(\|\varphi\|\le\varphi(1)\)，而在单位元处取值给出反向不等式。
\end{proof}

例如 \(M_n(\mathbb C)\) 上的归一化迹 \(\tau(a)=\operatorname{tr}(a)/n\) 是态，因为 \(\operatorname{tr}(a^*a)=\sum_{i,j}|a_{ij}|^2\)。对单位向量 \(v\)，\(\omega_v(a)=\langle v,av\rangle\) 也是态；\(C(X)\) 中的点取值则既是态又保持乘法。一般态不保持乘法，如 \(M_2(\mathbb C)\) 中 \(\tau(E_{11})^2=1/4\ne1/2=\tau(E_{11}^2)\)。

\ProseParagraphBreak
态却能提供长度和内积的候选：\(\varphi(a^*b)\) 对第一变量共轭线性、对第二变量线性，且平方值非负。这个型可能退化。例如取 \(v=e_1\)，则 \(\omega_v(a^*a)=\|ae_1\|^2\)；非零矩阵 \(E_{22}\) 的值为零，因为 \(E_{22}e_1=0\)。所以要从态得到真正的内积空间，必须先把所有零长度元素模去，再作完备化。下面的 GNS 构造正是将这两步与代数的左乘作用结合起来。

\subsection{GNS 构造与循环表示}
\label{b2:subsec:D03:02}

\begin{theorem}[Gelfand--Naimark--Segal 构造]\label{b2:thm:gns-unital}
设 \(A\ne0\) 为含幺复 \(C^*\)-代数，\(\varphi\) 为 \(A\) 上的态，则存在复 Hilbert 空间 \(H_\varphi\)、保幺 \(*\)-表示 \(\pi_\varphi:A\to B(H_\varphi)\) 和单位向量 \(\xi_\varphi\)，使对每个 \(a\in A\)，
\[
\varphi(a)=\bigl\langle\xi_\varphi,\pi_\varphi(a)\xi_\varphi\bigr\rangle,
\qquad\overline{\pi_\varphi(A)\xi_\varphi}=H_\varphi.
\]
这个三元组在保持指定向量并交织表示的酉同构下唯一。
\symidx{\(H_\varphi\)：态的 GNS Hilbert 空间}
\symidx{\(\pi_\varphi\)：态的 GNS 表示}
\symidx{\(\xi_\varphi\)：GNS 表示的指定循环向量}
\end{theorem}
\begin{proof}
令
\[
N_\varphi=\bigl\{\,a\in A:\varphi(a^*a)=0\,\bigr\}.
\]
\symidx{\(N_\varphi\)：GNS 构造中半内积的零空间}
由\cref{b2:lem:positive-functional-cs}，其中每个元素与任意元素的半内积都为零，故 \(N_\varphi\) 为线性子空间，且
\[
\bigl\langle[a],[b]\bigr\rangle=\varphi(a^*b)
\]
在商 \(A/N_\varphi\) 上良定义并正定。\cref{b2:prop:cstar-positive-input}还给出
\begin{equation}\label{b2:eq:gns-bounded-left-action}
\varphi\bigl((ca)^*(ca)\bigr)
=\varphi(a^*c^*ca)\le\|c\|^2\cdot\varphi(a^*a).
\end{equation}
确实，对正元素 \(\|c\|^2 1-c^*c\) 作 \(a^*(\cdot)a\) 并应用 \(\varphi\)，就得到此不等式。它同时证明 \(N_\varphi\) 为左理想，以及左乘 \([a]\mapsto[ca]\) 为范数至多 \(\|c\|\) 的有界算子。

把内积空间完备化为 \(H_\varphi\)，左乘唯一连续延拓为 \(\pi_\varphi(c)\)。乘法及单位元关系先在稠密商空间上成立，连续性使其在完备化后仍成立。又有
\[
\bigl\langle[ca],[b]\bigr\rangle=\varphi(a^*c^*b)
=\bigl\langle[a],[c^*b]\bigr\rangle,
\]
所以 \(\pi_\varphi(c)^*=\pi_\varphi(c^*)\)。取 \(\xi_\varphi=[1]\)，它的范数为一，\(\pi_\varphi(a)\xi_\varphi=[a]\)，因而满足循环性及泛函公式。

若 \((\rho,H,\eta)\) 是另一三元组，在稠密子空间上定义 \([a]\mapsto\rho(a)\eta\)。其内积为
\[
\bigl\langle\rho(a)\eta,\rho(b)\eta\bigr\rangle=\varphi(a^*b),
\]
故良定义且等距。循环性使其像稠密，所以延拓为满酉算子；它交织左作用并送指定向量到 \(\eta\)。这两个要求又在稠密向量 \([a]\) 上决定它，故唯一。
\end{proof}

这称为\term[GNS-gouzao]{GNS 构造}[Gelfand--Naimark--Segal construction]。\footnote{奈马克（\OriginalName{Марк Аронович Наймарк}，1909-1978），苏联数学家，研究泛函分析与算子表示。西格尔（Irving Ezra Segal，1918-1998），美国数学家，研究泛函分析与量子场论。}上面的完备化是必需的一步，不能把一般的代数商空间直接当作 Hilbert 空间。非含幺版本及其循环向量的构造可参阅\cite[\nopp II.6.4.1--II.6.4.3]{Blackadar2017OperatorAlgebrasCorrected}。

\ProseParagraphBreak
对 \(M_n(\mathbb C)\) 的迹态，\(N_\tau=0\)，所以 \(H_\tau=M_n(\mathbb C)\)，内积为 \(\operatorname{tr}(a^*b)/n\)，表示为左乘，循环向量为 \(I_n\)。维数为 \(n^2\)，并非自然列模的 \(n\)。对向量态 \(\omega_v\)，则 \(N_{\omega_v}=\{\,a:av=0\,\}\)，映射 \([a]\mapsto av\) 等距满射到 \(\mathbb C^n\)，恢复自然表示。不同态因此可以给出不同的循环模型。

\subsection{Gelfand--Naimark 表示定理}
\label{b2:subsec:D03:03}

一次 GNS 构造未必忠实。例如 \(C(X)\) 的点取值只看到一个点，所有在该点为零的函数都被零化。为了看见整个代数，需要足够多的态。

\begin{proposition}[达到范数的态]\label{b2:prop:norming-state-external}
设 \(A\ne0\) 为含幺复 \(C^*\)-代数。对每个正元素 \(p\in A_+\)，存在 \(A\) 上的态 \(\varphi\) 使 \(\varphi(p)=\|p\|\)。
\end{proposition}

\begin{proof}
由函数演算，\(\|p\|\) 属于 \(p\) 的谱；在 \(C^*(1,p)\cong C(\sigma_A(p))\) 上取该点的取值泛函 \(\psi\)。它满足 \(\|\psi\|=\psi(1)=1\)、\(\psi(p)=\|p\|\)。\term[Hahn-Banach-dingli]{Hahn--Banach 定理}[Hahn--Banach theorem]将它保范延拓为 \(A\) 上的复线性泛函 \(\varphi\)，仍有 \(\|\varphi\|=\varphi(1)=1\)。

核对延拓的正性。对自伴元 \(h\) 及实数 \(t\)，\(C^*\) 恒等式给出
\[
|1+it\cdot\varphi(h)|^2\le\|1+ith\|^2
=\|1+t^2h^2\|\le1+t^2\|h\|^2.
\]
左端关于 \(t\) 的一次项为 \(-2t\operatorname{Im}\varphi(h)\)。让 \(t\) 从正、负两侧趋于零，必得 \(\operatorname{Im}\varphi(h)=0\)。若 \(q\ge0\)、\(\|q\|\le1\)，正性的谱论基础使 \(\|1-q\|\le1\)，故 \(|1-\varphi(q)|\le1\)。因为 \(\varphi(q)\) 为实数，必有 \(\varphi(q)\ge0\)。任意正元素可缩放到这种情形，所以 \(\varphi\) 是态，并保留所需取值。
\end{proof}

这里仅用 Hahn--Banach 的保范延拓作为分析学基础；延拓何以仍为正已在证明中核对。这个存在性论证亦见\cite[\nopp II.6.3.1; Proposition II.6.3.3]{Blackadar2017OperatorAlgebrasCorrected}。

\begin{theorem}[Gelfand--Naimark，含幺情形]\label{b2:thm:gelfand-naimark-unital}
设 \(A\ne0\) 为含幺复 \(C^*\)-代数，则存在复 Hilbert 空间 \(H\) 及保幺等距 \(*\)-同态 \(A\rightarrow B(H)\)，其像为 \(B(H)\) 的含幺范数闭子代数。
\end{theorem}
\begin{proof}
记全部态组成的集合为 \(\mathcal S(A)\)。先取有限支集的代数直和 \(\bigoplus_{\varphi\in\mathcal S(A)}H_\varphi\)，以各分量内积之和为内积，再作完备化，得到 Hilbert 直和 \(H\)。因为各 \(H_\varphi\) 已经完备，\(H\) 可识别为
\[
\begin{gathered}
H=\left\{\,(\xi_\varphi)_{\varphi\in\mathcal S(A)}
\,\Bigm\vert\,
\sum_{\varphi\in\mathcal S(A)}\|\xi_\varphi\|^2<\infty\,\right\},
\\
\|\xi\|^2=\sum_\varphi\|\xi_\varphi\|^2.
\end{gathered}
\]
任意指标集上的平方和指有限子集部分和的上确界。因此这里允许平方可和的分量族，而不仅是有限支集向量。

令 \(\pi(a)(\xi_\varphi)_\varphi=(\pi_\varphi(a)\xi_\varphi)_\varphi\)。式\eqref{b2:eq:gns-bounded-left-action}给出与态无关的统一界 \(\|\pi_\varphi(a)\|\le\|a\|\)，于是
\[
\sum_\varphi\|\pi_\varphi(a)\xi_\varphi\|^2
\le\|a\|^2\sum_\varphi\|\xi_\varphi\|^2.
\]
故 \(\pi(a)\) 确实把 \(H\) 送入自身，且 \(\|\pi(a)\|\le\|a\|\)。乘法、单位与对合关系先在有限支集的稠密子空间上成立，再由连续性延伸到 \(H\)，所以 \(\pi\) 是含幺 \(*\)-同态。对给定 \(a\)，应用\cref{b2:prop:norming-state-external}于 \(p=a^*a\)，选择相应态，把其单位循环向量 \(\xi_\varphi\) 放在 \(H\) 的该分量、其他分量置零，则
\[
\bigl\lVert\pi(a)\bigr\rVert^2\ge\bigl\lVert\pi_\varphi(a)\xi_\varphi\bigr\rVert^2
=\varphi(a^*a)=\|a^*a\|=\|a\|^2.
\]
结合相反方向的统一界得到等距性，特别是单射。等距像因 \(A\) 完备而闭，结论成立。
\end{proof}

这就是含幺情形的\term[Gelfand-Naimark-biaoshi-dingli]{Gelfand--Naimark 表示定理}[Gelfand--Naimark representation theorem]。本证明先构造态，再以它们的 GNS 表示组成忠实表示。

\begin{proposition}[添幺与非含幺表示]\label{b2:prop:cstar-unitization-representation}
设 \(A\ne0\) 为不含幺的复 \(C^*\)-代数。它可等距嵌入一个含幺 \(C^*\)-代数 \(\widetilde A\)，作为闭双边 \(*\)-理想，且 \(\widetilde A/A\cong\mathbb C\)。特别地，\(A\) 有忠实等距 Hilbert 空间表示。
\end{proposition}
\begin{proof}
取向量空间 \(\widetilde A=A\oplus\mathbb C1\)，规定
\[
(a+\lambda1)(b+\mu1)=ab+\lambda b+\mu a+\lambda\mu1,
\qquad (a+\lambda1)^*=a^*+\bar\lambda1.
\]
把它作用于 Banach 空间 \(A\)，令 \(L_{a+\lambda1}x=ax+\lambda x\)，定义 \(N(c)=\|L_c\|\)。若 \(N(a+\lambda1)=0\) 且 \(\lambda\ne0\)，\(-a/\lambda\) 就是 \(A\) 的左单位；取对合得到右单位，两者相同，违反不含幺假设。若 \(\lambda=0\)，则 \(aa^*=0\)，故 \(a=0\)。因此 \(N\) 是范数，且由 \(x=a^*/\|a\|\) 得 \(N(a)=\|a\|\)。

对 \(c\in\widetilde A\)、\(x\in A\)，有
\[
\|c^*x\|^2=\|x^*cc^*x\|
\le\|x\|\cdot N(c)\cdot\|c^*x\|.
\]
消去非零因子并取上确界，得 \(N(c^*)\le N(c)\)；取对合又得相等。再由
\[
\|cx\|^2=\|x^*c^*cx\|\le\|x\|^2\cdot N(c^*c)
\]
和次乘性，得到 \(N(c)^2=N(c^*c)\)。

还须检查完备性。必有 \(|\lambda|\le N(a+\lambda1)\)：否则在 \(B(A)\) 中有 \(\|I+L_{a/\lambda}\|<1\)。写成
\[
-L_{a/\lambda}=I-\bigl(I+L_{a/\lambda}\bigr),
\]
右侧的括号内算子范数小于一，故 Neumann 级数给出 \(-L_{a/\lambda}\) 的逆，从而 \(L_a\) 可逆。取 \(e\in A\) 满足 \(ae=a\)，则 \(a(ex-x)=0\) 及 \(L_a\) 单射给出 \(ex=x\)，又会产生单位元。因此 \(\|a\|\le2N(a+\lambda1)\)。一个 \(N\)-柯西列的 \(A\) 坐标和复数坐标分别柯西，取两坐标的极限便得它的 \(N\)-极限。于是 \(\widetilde A\) 确为含幺 \(C^*\)-代数，\(A\) 为其中闭理想，取复数坐标给出商同构。将上一定理的等距表示限制到 \(A\)，即得非含幺结论。
\end{proof}

这个表示的原代数像未必含有 Hilbert 空间的单位算子；添幺和一般表示的比较见\cite[\nopp II.1.2; Corollary II.6.4.10]{Blackadar2017OperatorAlgebrasCorrected}。

\begin{proposition}[一般 \(C^*\)-代数的 \(*\)-同态]\label{b2:prop:cstar-star-hom-general}
设 \(A,B\) 为复 \(C^*\)-代数，\(\pi:A\to B\) 为 \(*\)-同态，不要求代数含幺或映射保幺。则对每个 \(a\in A\)，有
\[
\|\pi(a)\|\le\|a\|.
\]
因此 \(\pi\) 自动连续。若 \(\pi\) 单射，则它等距，即 \(\|\pi(a)\|=\|a\|\)。
\end{proposition}
\begin{proof}
零代数的情形直接成立。若 \(A\) 含幺，结论由\cref{b2:prop:cstar-star-hom-contractivity-unital,b2:prop:cstar-star-hom-isometry-unital}给出。

设 \(A\) 不含幺，取它的等距添幺 \(\widetilde A\)。若 \(B\) 含幺，令 \(B^\sharp=B\)；否则取 \(B^\sharp=\widetilde B\)。公式
\[
\widetilde\pi(a+\lambda1)=\pi(a)+\lambda1_{B^\sharp}
\]
定义保幺 \(*\)-同态 \(\widetilde\pi:\widetilde A\to B^\sharp\)。含幺情形的收缩性及添幺的等距嵌入给出
\(\|\pi(a)\|\le\|a\|\)。

再设 \(\pi\) 单射。若 \(\widetilde\pi(a+\lambda1)=0\)，当 \(\lambda=0\) 时已有 \(a=0\)；若 \(\lambda\ne0\)，则 \(e=-a/\lambda\in A\) 满足 \(\pi(e)=1_{B^\sharp}\)。对每个 \(x\in A\)，于是
\[
\pi(ex)=\pi(x)=\pi(xe).
\]
由单射性得 \(ex=xe=x\)，这与 \(A\) 不含幺矛盾。因此 \(\widetilde\pi\) 也单射，含幺情形的等距性限制到 \(A\) 即得所需结论。
\end{proof}
